\section*{Chapter \ref{ch:ml:trees-and-boosting} --- Trees and Boosting}

\begin{solution}{exo:ml:trees-and-boosting:1}
$0.3\sigma^2 + 0.7\sigma^2/50 = 0.314\sigma^2$; with infinitely many trees, $0.3\sigma^2$: only decorrelation lowers the
floor.
\end{solution}

\begin{solution}{exo:ml:trees-and-boosting:2}
$400\times(-0.1 - 0.08)^2 + 600\times(0.2 - 0.08)^2 = 12.96 + 8.64 = 21.6$.
\end{solution}

\begin{solution}{exo:ml:trees-and-boosting:3}
$F_1 = 0.1\times0.8 = 0.08$, residual 0.92. If each tree predicted the current residual exactly, each step would keep
$1 - \eta = 0.9$ of it: $0.9^{10} = 0.349$ after ten steps.
\end{solution}

\begin{solution}{exo:ml:trees-and-boosting:4}
An out-of-bag row shares its month, hence its factor returns, with hundreds of in-bag rows of similar characteristics;
the trees have learned that month's factor moves from them, so the row is not out of sample. The test months share
nothing with the training months.
\end{solution}

\begin{solution}{exo:ml:trees-and-boosting:5}
The defaults average 0.09\% with a standard deviation of 0.07 points across seeds; 0.14\% on one seed is within one
standard deviation of their mean. Rescore the candidate with the same four seeds (and on purged folds) before calling it
better.
\end{solution}

\begin{solution}{exo:ml:trees-and-boosting:6}
The number of trees was chosen on the test months, so the test score is a selection score, optimistic by
\cref{prop:ml:validation:max}. Stop on a validation block purged from both the training and the test months.
\end{solution}

\begin{solution}{exo:ml:trees-and-boosting:7}
\texttt{monotone(ofi5\_sign=-1)}: R-squared 0.305, against 0.333 unconstrained and 0.364 with the right sign; raising
the imbalance now lowers the forecast for 95\% of the rows. A wrong constraint is a wrong prior imposed exactly; it costs
more than no constraint.
\end{solution}

\begin{solution}{exo:ml:trees-and-boosting:8}
With squared loss and exact fits, $F_m = F_{m-1} + \eta(y - F_{m-1})$, so $y - F_m = (1 - \eta)(y - F_{m-1})$ and, from
$F_0 = 0$, $y - F_M = (1 - \eta)^My$. The fit's progress depends on $\eta M$ roughly (for small $\eta$, $(1 - \eta)^M
\approx e^{-\eta M}$): halving the \omterm{def:ml:trees-and-boosting:boosting}{learning rate} doubles the trees needed, and stopping early at a given $\eta$ is like
choosing a smaller effective $\eta M$. Real trees fit residuals only partly and fit noise too, which is why a small
$\eta$ with more trees generalises better than a large one with few.
\end{solution}

\begin{solution}{pb:ml:trees-and-boosting:1}
\begin{enumerate}
\item 0.24\%, $-0.20\%$, $-3.42\%$: deeper trees fit more noise (the variance term).
\item 0.20\% out of sample, 0.42\% out of bag, 2.76\% in sample.
\item Out-of-bag rows share their month's factor returns with in-bag rows.
\item More trees remove the $(1-\rho)\sigma^2/B$ term; only less correlated trees lower $\rho\sigma^2$.
\item $-0.19\%$.
\item After 34 trees at 0.1 (test 0.05\%) and 233 trees at 0.01 (test 0.26\%).
\item At 0.1, the test R-squared falls to $-3.38\%$; at 0.01, to 0.03\%.
\item Each tree moves the fit a little, so noise must be fitted slowly and the validation curve is flat near its peak.
\item From $-0.54\%$ to 0.39\%.
\item Four leaves, \omterm{def:ml:trees-and-boosting:boosting}{learning rate} 0.02, 200 trees, 20 rows per leaf; 0.36--0.39\% with four seeds, 0.32\% on the test.
\item Standard deviations of 0.07 points for the defaults, 0.014 for the winner: five times noisier.
\item Four leaves, 1\,000 rows per leaf, no subsampling: 0.35\% on validation, 0.31\% on test.
\item R-squared 0.333 to 0.364; violations from 7.0\% of the rows to none.
\item It removes directions of fit that only noise supports: a restriction that the truth satisfies lowers variance
without adding bias.
\item \emph{Named result}: the defaults' validation score varies with the seed by 0.07 points (0.00--0.17\%) and the
winner's by 0.014, against a gain from tuning of about 0.3 points (0.09\% to 0.39\%); the winner loses 0.07 points from
validation to test.
\item The plateau choice: as good on test, less flexible, and less sensitive to the seed.
\item Three to five, enough to see a standard deviation; more if the candidates are within one of each other.
\item Every configuration and its score, the seeds, the validation block, and the rule used to choose.
\item The effective sample of the characteristics' effects is the months, so more stocks help less than more months;
complex configurations would still be noisy.
\item Small \omterm{def:ml:trees-and-boosting:boosting}{learning rates}, few leaves, many rows per leaf, \omterm{def:ml:trees-and-boosting:early}{early stopping} on purged data, several seeds, and a choice
from the plateau.
\end{enumerate}
\end{solution}

\begin{solution}{iq:ml:trees-and-boosting:1}
\omterm{def:ml:trees-and-boosting:bagging}{Bagging} averages models fitted independently to bootstrap samples, reducing variance; boosting fits models in sequence,
each to the errors of the ensemble so far, reducing bias (and, with shrinkage and \omterm{def:ml:trees-and-boosting:early}{early stopping}, controlling variance).
\iqlookfor{parallel versus sequential, variance versus bias.}
\end{solution}

\begin{solution}{iq:ml:trees-and-boosting:2}
The ones that control variance: a small \omterm{def:ml:trees-and-boosting:boosting}{learning rate} with \omterm{def:ml:trees-and-boosting:early}{early stopping}, few leaves or shallow depth, a large minimum
number of rows per leaf, row and column subsampling, an $\ell_2$ penalty. All in the direction of less flexibility.
\iqlookfor{regularisers first, and the direction on low signal-to-noise data.}
\end{solution}

\begin{solution}{iq:ml:trees-and-boosting:3}
Out-of-bag rows share dates (and common factor shocks, overlapping labels) with in-bag rows, so they are not independent
test data; the score is optimistic. Use purged, time-ordered validation.
\iqlookfor{the dependence between rows of a panel.}
\end{solution}

\begin{solution}{iq:ml:trees-and-boosting:4}
A restriction that the model be non-decreasing (or non-increasing) in a feature. Candidates: order-flow imbalance and
queue imbalance (a higher bid-side pressure should not predict a lower price); a borrow fee in a short-squeeze model; a
distance to a price limit.
\iqlookfor{the definition and economically justified signs.}
\end{solution}

\begin{solution}{iq:ml:trees-and-boosting:5}
Compare with the noise first: rescore both with several seeds and on purged folds; if they are within one standard
deviation, take the simpler, less turnover-prone, or better-constrained one.
\iqlookfor{noise before ranking, then simplicity.}
\end{solution}

\begin{solution}{iq:ml:trees-and-boosting:6}
The negative gradient of $\frac12(y - F)^2$ in $F$ is $y - F$, the residual, so each tree fits residuals. With the Huber
loss the gradient is the residual clipped at $\pm\delta$: large errors (fat-tailed returns) pull the fit less.
\iqlookfor{the gradient computation and the robustness of the Huber loss.}
\end{solution}
